\documentclass[11pt,a4paper]{article}
\usepackage[utf8]{inputenc}
\usepackage[vietnamese]{babel}
\usepackage{amsmath,amssymb}
\usepackage{geometry}
\geometry{a4paper, total={180mm,260mm}, left=15mm, right=15mm, top=18mm, bottom=18mm}

\usepackage{tikz}
\usetikzlibrary{shapes.geometric, arrows.meta, positioning, calc}
\tikzset{
  arr/.style={-{Stealth[scale=0.8]}},
  every node/.append style={align=center}
}
\usepackage{tcolorbox}
\tcbuselibrary{skins,breakable}
\usepackage{fontawesome}
\usepackage{fancyhdr}
\usepackage{enumitem}

\pagestyle{fancy}
\fancyhf{}
\renewcommand{\headrulewidth}{0.5pt}
\renewcommand{\footrulewidth}{0.5pt}
\fancyhead[L]{\small\bfseries\color{violet!85!black}\faBook\ SỔ TAY KIẾN THỨC TOÁN 2}
\fancyhead[R]{\small\bfseries\color{violet!85!black}Lớp Toán Thầy Lĩnh -- loptoanthaylinh.vn}
\fancyfoot[C]{\small\color{gray}-- Sổ tay kiến thức: Bài 18. Luyện tập chung Chủ đề 3 --}

\setlength{\parindent}{0pt}
\setlength{\parskip}{0pt}

\begin{document}

% --- BANNER BÀI HỌC ---
\begin{tcolorbox}[
  enhanced,
  colback=violet!8!white,
  colframe=violet!70!black,
  arc=3mm, boxrule=1.2pt,
  left=3mm, right=3mm, top=2mm, bottom=2mm,
  halign=center
]
  {\LARGE\bfseries\color{violet!90!black}\faBookmark\ BÀI 18: LUYỆN TẬP CHUNG CHỦ ĐỀ 3}\\[1mm]
  {\itshape\footnotesize\color{black!75} Hệ thống hóa toàn bộ kiến thức và bài toán có lời văn về Ki-lô-gam và Lít trên 1 trang A4}
\end{tcolorbox}

\vspace{-2mm}

\noindent
\begin{minipage}[t]{0.49\textwidth}\vspace{0pt}
\begin{tcolorbox}[
  enhanced,
  colback=violet!4!white,
  colframe=violet!75!black,
  arc=2mm, boxrule=0.8pt,
  left=2mm, right=2mm, top=1.8mm, bottom=1.8mm,
  title={\footnotesize\bfseries\color{white}\faKey\ 1. Hai đại lượng thực hành}
]
{\footnotesize
\begin{itemize}\setlength{\itemsep}{1pt}\setlength{\leftskip}{-2.5mm}
  \item \textbf{Khối lượng:} Đo bằng \textbf{ki-lô-gam (kg)}.
  \item \textbf{Dung tích:} Đo bằng \textbf{lít ($l$)}.
  \item \textbf{Quy tắc tính toán:}\\
  Thực hiện cộng, trừ bình thường, giữ nguyên tên đơn vị đo ở kết quả.
  \item \textbf{Quy tắc so sánh:}\\
  Cùng đơn vị đo thì so sánh số: $18\text{ kg} < 20\text{ kg}$; $45\ l > 40\ l$.
\end{itemize}
}
\end{tcolorbox}
\end{minipage}
\hfill
\begin{minipage}[t]{0.49\textwidth}\vspace{0pt}
\begin{tcolorbox}[
  enhanced,
  colback=teal!4!white,
  colframe=teal!75!black,
  arc=2mm, boxrule=0.8pt,
  left=2mm, right=2mm, top=1.8mm, bottom=1.8mm,
  title={\footnotesize\bfseries\color{white}\faSliders\ 2. Kĩ năng giải toán lời văn}
]
{\footnotesize
\begin{itemize}\setlength{\itemsep}{1pt}\setlength{\leftskip}{-2.5mm}
  \item \textbf{Bài toán tìm tổng:} Gộp lại, cả hai thùng, cả hai bao $\implies$ Phép cộng ($+$).
  \item \textbf{Bài toán tìm còn lại:} Bán đi, rót bớt ra, dùng hết $\implies$ Phép trừ ($-$).
  \item \textbf{Bài toán nhiều hơn/ít hơn:}
  \begin{itemize}\setlength{\itemsep}{0.5pt}\setlength{\leftskip}{-2.5mm}
    \item Tìm bao to nặng hơn $\implies$ Lấy số bé cộng phần hơn.
    \item Tìm thùng nhỏ chứa ít hơn $\implies$ Lấy số lớn trừ phần kém.
  \end{itemize}
\end{itemize}
}
\end{tcolorbox}
\end{minipage}

\vspace{-1.5mm}
\begin{tcolorbox}[
  enhanced,
  colback=orange!4!white,
  colframe=orange!80!black,
  arc=2mm, boxrule=0.8pt,
  left=2mm, right=2mm, top=1.5mm, bottom=1.5mm,
  title={\footnotesize\bfseries\color{white}\faSitemap\ 3. Sơ đồ tư duy: Cặp bài trùng Đại lượng đo lường lớp 2}
]
\centering
\begin{tikzpicture}[x=1.5cm, y=0.8cm]
  \node[rectangle, rounded corners=2pt, draw=violet!80!black, fill=violet!20, font=\scriptsize\bfseries] (center) at (0,0) {ĐẠI LƯỢNG\\LỚP 2};
  \node[rectangle, rounded corners=2pt, draw=teal!80!black, fill=teal!15, font=\scriptsize\bfseries, left=1.2cm of center] (kg) {KI-LÔ-GAM (kg)\\Cân nặng của vật};
  \node[rectangle, rounded corners=2pt, draw=orange!80!black, fill=orange!15, font=\scriptsize\bfseries, right=1.2cm of center] (lit) {LÍT ($l$)\\Sức chứa chất lỏng};

  \draw[arr, thick, teal!80!black, <-] (kg) -- (center);
  \draw[arr, thick, orange!80!black, ->] (center) -- (lit);
\end{tikzpicture}
\end{tcolorbox}

\vspace{-1.5mm}
\noindent
\begin{minipage}[t]{0.49\textwidth}\vspace{0pt}
\begin{tcolorbox}[
  enhanced,
  colback=green!4!white,
  colframe=green!60!black,
  arc=2mm, boxrule=0.8pt,
  left=2mm, right=2mm, top=1.8mm, bottom=1.8mm,
  title={\footnotesize\bfseries\color{white}\faCheckCircle\ 4. Ví dụ thực hành chuẩn mực}
]
{\footnotesize
\textbf{Bài toán tổng hợp:}\\
Can thứ nhất có $18\ l$ mật ong, can thứ hai có ít hơn can thứ nhất $5\ l$ mật ong. Hỏi can thứ hai có bao nhiêu lít mật ong?\\
\textit{Bài giải:}\\
Can thứ hai có số lít mật ong là:
$$18 - 5 = 13\text{ (}l\text{)}$$
Đáp số: $13\ l$ mật ong.
}
\end{tcolorbox}
\end{minipage}
\hfill
\begin{minipage}[t]{0.49\textwidth}\vspace{0pt}
\begin{tcolorbox}[
  enhanced,
  colback=red!4!white,
  colframe=red!75!black,
  arc=2mm, boxrule=0.8pt,
  left=2mm, right=2mm, top=1.8mm, bottom=1.8mm,
  title={\footnotesize\bfseries\color{white}\faExclamationTriangle\ 5. Mẹo tránh bẫy trắc nghiệm}
]
{\footnotesize
\begin{itemize}\setlength{\itemsep}{1pt}\setlength{\leftskip}{-2.5mm}
  \item \textbf{Cẩn thận đơn vị trong lời giải:}\\
  Lời giải hỏi "bao nhiêu lít" thì đơn vị phép tính là ($l$). Hỏi "bao nhiêu ki-lô-gam" thì đơn vị là (kg).
  \item Khi tính nhẩm dãy phép tính, phải ghi đơn vị ở kết quả cuối cùng: $12\text{ kg} + 8\text{ kg} - 5\text{ kg} = 15\text{ kg}$.
\end{itemize}
}
\end{tcolorbox}
\end{minipage}


\end{document}
